\documentclass[10pt]{article}
\usepackage[margin=20mm]{geometry}
\usepackage{amsmath,amssymb,mathtools,fontspec,unicode-math,xurl,hyperref}
\IfFontExistsTF{DejaVu Serif}{\setmainfont{DejaVu Serif}}{\setmainfont{Latin Modern Roman}}
\setmathfont{latinmodern-math.otf}
\setlength{\parindent}{0pt}\setlength{\parskip}{6pt}\sloppy
\title{A Liouville theorem in the Heisenberg group: normalization and boundary terms}
\author{AI-assisted technical review}
\date{}
\begin{document}
\maketitle
Abstract: Two Major issues affect the proof of the higher-dimensional classification in the bound author manuscript. Its stated Bochner inequality fails on an explicit positive entire solution with the prescribed normalization. The subsequent annular integration by parts omits an inner-boundary flux; a translated solution makes the omitted contribution strictly positive in the upper estimate. Two local normalization errors are also recorded. These findings do not disprove either classification conclusion. The weighted argument in homogeneous dimension four was checked separately, with the stated limitations concerning its external rigidity input.\par
\section*{Source, assumptions, and scope}
This report concerns the 23-page CVGMT-hosted manuscript by Giovanni Catino, Yanyan Li, Dario D. Monticelli, and Alberto Roncoroni, A Liouville theorem in the Heisenberg group [paper]. The exact source has 438267 bytes and SHA256 \[\text{\textnormal{
7d4636e5a79dcc6e4d8a0f4a57b547417de162a7b1ca73d4f772cc8bffb2cb05}}.\] All 23 pages, including references, were read as extracted text. CVGMT labels its record an accepted paper and gives 27 November 2024 as the record’s update date. That date is not an authenticated revision date for this PDF. Exact alignment with the accepted version, with arXiv:2310.10469, or with the journal version was not established. The associated journal DOI is 10.4171/JEMS/1705; the publisher records online publication on 15 September 2025 [journal]. Findings below are attributed only to the bound CVGMT file.\par
The equation is \[-\Delta_{\mathbb H^n}u=2n^2u^{1+2/n},\qquad u>0\quad\hbox{on }\mathbb H^n.\] Theorem 1.1 treats \(n=1\) without a finite-energy or symmetry assumption. Theorem 1.2 treats \(n\geq2\) under \(u(\xi)\leq C/(1+|\xi|^n)\). This borderline decay alone need not give finite critical energy. No unconditional classification in every dimension is attributed to Theorem 1.2.\par
In the manuscript’s convention, \[X_j=\partial_{x_j}+2y_j\partial_t,\quad
 Y_j=\partial_{y_j}-2x_j\partial_t,\quad
 \Delta_{\mathbb H^n}=\tfrac12\sum_j(X_j^2+Y_j^2),\quad
 |\nabla f|^2=\tfrac12\sum_j((X_jf)^2+(Y_jf)^2).\] Thus \(|\nabla f|^2=2|\partial f|^2\). Writing \(f=n^{-1}\log u\) gives \[\Delta_{\mathbb H^n}f=-n|\nabla f|^2-2ne^{2f},\qquad
 g=|\partial f|^2+e^{2f}-if_0,\quad f_0=\partial_tf.\] The factors of two matter in both findings below.\par
\section*{A false Bochner input}
Equation reference: badbochner. Problem 1 (Major: equation (4.4), page 14).  The displayed inequality, stated for every positive function \(\nu\), is \[
 \Delta_{\mathbb H^n}|\nabla f|^2\geq
 \frac{(\Delta_{\mathbb H^n}f)^2}{n}+nf_0^2+
 2\langle\nabla f,\nabla\Delta_{\mathbb H^n}f\rangle
 -\frac2\nu|\nabla f|^2-2\nu|\nabla f_0|^2.\] It is false even when \(f\) is obtained from a positive solution satisfying the assumptions of Proposition 4.1.\par
Take \(n=2\), put \(\rho=|z|^2\) and \(S=(1+\rho)^2+t^2\), and set \[u=S^{-1},\qquad f=-\tfrac12\log S.\] Direct differentiation gives \(-\Delta_{\mathbb H^2}u=8u^2\) identically. The function is smooth and positive on the entire group and satisfies \(u\leq1\) everywhere. In particular, the local upper bound in Proposition 4.1 holds with center zero, \(R=1\), and \(a=1\).\par
At \(z=(1/2,0)\), \(t=0\), and \(\nu=5/4\), the relevant exact quantities are \[\begin{array}{c|c}
\text{quantity}&\text{value}\\ \hline
|\nabla f|^2&8/25\\
|\nabla f_0|^2&128/625\\
\Delta_{\mathbb H^2}f&-16/5\\
f_0&0\\
2\langle\nabla f,\nabla\Delta_{\mathbb H^2}f\rangle&-256/125\\
\Delta_{\mathbb H^2}|\nabla f|^2&896/625
\end{array}\] The right side of [badbochner] (\#badbochner) is \(256/125=1280/625\). Consequently its left side minus its right side is \(-384/625\).\par
These computations were independently executed symbolically with rational arithmetic. One implementation used, for a radial function \(h(\rho,t)\), \[\Delta_{\mathbb H^n}h=2\rho(h_{\rho\rho}+h_{tt})+2nh_\rho,\quad
 \langle\nabla f,\nabla h\rangle
 =2\rho(f_\rho h_\rho+f_th_t).\] A second implementation differentiated the real fields \(X_j,Y_j\) directly in all four horizontal variables. Both gave the displayed values, and the equation for \(u\) was checked as an identity, not just at the selected point. Thus this is not a test function that fails the governing differential equation.\par
Equation (4.4) is used to obtain the differential inequality for the maximum-principle quantity in Proposition 4.1. That proposition supplies the gradient decay required in Theorem 1.2. The defect therefore affects a core proof dependency. A repaired proof must derive a Bochner inequality with constants consistent with these vector fields and then redo the torsion absorption and barrier estimates. No complete replacement argument is claimed here. The example is itself a classified bubble, so it does not contradict the proposed classification.\par
\section*{The inner boundary of the annulus}
Problem 2 (Major: equation (4.9) through (4.10), page 18).  The proof of Theorem 1.2 integrates by parts over \(A_R=B_R\setminus B_{R/2}\) using a cutoff \(\eta\) equal to one on \(B_{R/2}\). The multiplier \(\eta^{s-2}\) therefore does not vanish on the inner boundary. The displayed integration-by-parts equality omits that boundary contribution.\par
Write \(W=e^{2(n-1)f}\), and define the complex components \(V_\alpha=W f_\alpha\overline g\). Since the real vector fields underlying \(Z_\alpha\) and \(Z_{\bar\alpha}\) are divergence free for Haar measure, the relevant product rule is \[\operatorname{Re}\sum_\alpha Z_{\bar\alpha}
       (\eta^{s-2}V_\alpha)
 =W\eta^{s-2}\operatorname{Re}
       \sum_\alpha(f_{\alpha\bar\alpha}\overline g)
       +\text{the three interior product-rule terms}.\] The negative integral of the second-derivative term is the sum of those interior terms minus the total boundary flux. The outer flux vanishes because \(\eta=0\) there. No such argument removes the inner flux.\par
Equation reference: divfield. The missing term need not be zero or have a sign that permits its removal in an upper bound. For the \(n=2\) bubble of the preceding section, direct calculation gives \[
 \operatorname{Re}\sum_\alpha Z_{\bar\alpha}V_\alpha
   =\frac{2(\rho^2-1-t^2)}{((1+\rho)^2+t^2)^3}.\] This identity was checked with the actual complex vector fields. Translate this bubble by the group element \((z_0,0)\) with \(z_0=(10,0)\), and take \(R=2\). It remains a positive entire solution and satisfies the decay hypothesis of Theorem 1.2 for some finite constant. On the inner Koranyi ball \(B_1\) about zero, its coordinates relative to its center satisfy \[|z'|\geq9,\qquad |t'|\leq1+20=21.\] The numerator factor in [divfield] (\#divfield) is consequently bounded below by \[|z'|^4-1-(t')^2\geq9^4-1-21^2=6119>0.\] The integral of the divergence over \(B_1\) is strictly positive. The inner boundary has the opposite orientation as a boundary component of \(A_2\), so the annular inner flux is strictly negative. Its subtraction in the integration-by-parts formula adds a strictly positive term. The displayed equality on page 18 omits exactly this contribution.\par
This refutes the integration identity used to reach (4.10). It does not, by itself, prove that every possible inequality with an unspecified constant in (4.10) is false. The current proof supplies neither the omitted flux nor a bound for it. A possible repair is to introduce a cutoff vanishing at both boundaries of a larger annulus and keep every resulting derivative term; the required estimates and absorptions would then have to be proved. The compactly supported integration over the whole ball in Section 3 does not incur this particular boundary error. The present finding is independent of the Bochner coefficient problem.\par
\section*{Two local normalization corrections}
Problem 3 (Minor: the amplitude in equation (1.3), page 2).  For \[U_{\lambda,\mu}(z,t)=
 \frac{C}{|t+i|z|^2+z\cdot\mu+\lambda|^n},\qquad
 \operatorname{Im}\lambda>|\mu|^2/4,\] the amplitude is written as \(C=C(n,\lambda)\). It also depends on \(\mu\).\par
In the stated normalization the required amplitude is \[C=(\operatorname{Im}\lambda-|\mu|^2/4)^{n/2}.\] Indeed, complete the square with \(z'=z+i\overline\mu/2\) and \(t'=t+\operatorname{Re}(z\cdot\mu)+\operatorname{Re}\lambda\). The remaining scale is \(a=\operatorname{Im}\lambda-|\mu|^2/4>0\). For an explicit check, take \(n=1\), \(\lambda=i\), and evaluate at the origin. With \(\mu=0\), the equation requires \(C^2=1\); with \(\mu=1\), it requires \(C^2=3/4\). Both choices satisfy the positivity constraint, and the corresponding values of \(-\Delta_{\mathbb H^1}U(0)\) are respectively \(2C\) and \(3C/2\), while \(2U(0)^3=2C^3\) in each case. This corrects the parametrization without producing an extra solution outside the intended bubble family.\par
Problem 4 (Minor: equations (1.4)–(1.5), page 2).  The Euclidean background equation has coefficient \(n(n-2)/4\), whereas its displayed bubble has the normalization for coefficient \(n(n-2)\).\par
For \(n\geq3\) and \[V(x)=\left(\frac{\lambda}{\lambda^2+|x-x_0|^2}\right)^{(n-2)/2},\] one has \(-\Delta V=n(n-2)V^{(n+2)/(n-2)}\). Already \(n=3\), \(\lambda=1\), and \(x=x_0=0\) give \(-\Delta V(0)=3\), against the printed right side \(3/4\). Either multiply \(V\) by \(2^{(n-2)/2}\) or replace the coefficient in (1.4) by \(n(n-2)\). This is a local background normalization error; the Heisenberg proof does not rely on that displayed Euclidean parametrization.\par
\section*{Coverage of the remaining proof}
The definitions and contractions in Section 2 were followed with the manuscript’s derivative order and with repeated indices summed. The equation for \(f\), the relation \(f_{\alpha\bar\alpha}=-ng\), and the positivity \(|g|\geq e^{2f}\) were checked. The Jerison–Lee identity in Proposition 2.1 is used as an external differential identity; its long expansion and the rigidity conclusion from its vanishing were not independently reconstructed in full.\par
For \(n=1\), the weighted integral argument was checked separately from Section 4. Its weight \(\Psi=|g|^2e^{-2f}\) is positive. The vector-field and weight-derivative bounds in Lemma 2.2 give the stated absorption condition \(12\beta<1\). The unweighted Lemma 2.3 follows by the corresponding Cauchy–Schwarz absorption. The superharmonic lower bound in Proposition 2.4 is combined with positivity on compact sets, so negative powers of \(u\) have the polynomial bounds used on \(B_R\).\par
In Section 3 the restrictions on small positive \(\beta\) can be met simultaneously. For example, \(\beta=1/100\) and \(s=8\) give \[12\beta<1,\quad 2-36\beta>0,\quad 2-40\beta>0,\quad
 s-\frac{6-8\beta}{1-4\beta}=\frac{11}{6}>0.\] The estimates for \(I_3'\), the Young inequality in the estimate for \(I_4'\), and the absorption of \(I_2\) lead to the stated uniform weighted bound for \(I_1\) and the bound \(I_2\leq CR^2\). In particular the intermediate factor \(R^{36\beta-16\beta^2}\) is bounded by \(R^{36\beta}\) for \(R\geq1\). The final annular tail tends to zero because the weighted nonnegative density is integrable. No limit \(\beta\downarrow0\) and no unproved finite-energy assumption is needed at that step. These checks did not identify a further defect in that argument.\par
For the higher-dimensional proof, the parameter \(t\in[0,1]\) in Proposition 4.1 is a homotopy parameter, distinct from the group’s vertical coordinate. Conditional on its earlier differential inequality, the continuity of the maximum and threshold-crossing argument is coherent. Applying a local gradient estimate to the decay assumption requires balls away from the origin, which can be selected on the annulus; an origin-centered large ball does not satisfy a uniformly small upper bound merely from decay at infinity. The final scaling uses \(e^{2(n-1)f}\lesssim R^{-(Q-4)}\), \(|\partial f|^2\lesssim R^{-2}\), and volume of order \(R^Q\). Those exponents are compatible with the claimed annular estimates, but they do not repair the two Major issues established above.\par
\section*{Verification limits and disposition}
An independent Python/SymPy program executed the differential identities, the rational point evaluations, the positive flux lower bound, the two amplitude checks, and the parameter inequalities. It did not execute author code. Its first run completed the mathematical assertions but failed while serializing a SymPy Boolean value to JSON. That failed run is retained. A second run converted the output flags to native Boolean values and re-executed the checks successfully; both scripts, commands, exit codes, and outputs are retained.\par
The CVGMT and publisher records, the available preprint listing, and title-specific searches were checked for associated supplementary material, code, or corrections. No separate exact-version supplement or paper-specific code repository was located in this bounded search. Absence of code is not a mathematical defect. The official AMS route to the cited Jerison–Lee article returned HTTP 403, and its Section 3 rigidity proof was not read in full in this review. The classification conclusions are therefore not independently certified by this report.\par
The original image-read failure and its EMFILE records are preserved. In this same-chat continuation, all 23 pages of the exact bound source were actually displayed through the authorized file connection and inspected in the scientific session. The displayed conventions, quantifiers, coefficients in (4.4), and annular domains in (4.9)–(4.10) agree with the evidence above. This inspection did not change the scientific findings or remove the external-foundation limitations.\par
This revision is genuinely recompiled with the existing official Tectonic 0.17.0 and retained per-paper dependencies. Its delivery record binds the source, revised TeX, compiler, dependency payloads, resulting PDF, and actual observations of every resulting report page. The earlier false visual states and unsuccessful executions remain in the paper’s verification history.\par
The supported overall severity is Major, based on the two independent proof defects in Section 4. The two normalization corrections are Minor. These are initial-review findings for the exact bound source, not an independent acceptance decision. The scientific decision remains pending and the human-review list remains empty. Delivery completion is recorded only after the current source and report views, compilation, and required artifact bindings have actually passed their checks; it denotes completed initial-review delivery and is not independent scientific acceptance.\par
9 G. Catino, Y. Li, D. D. Monticelli, and A. Roncoroni, A Liouville theorem in the Heisenberg group, bound 23-page CVGMT author manuscript. https://cvgmt.sns.it/media/doc/paper/6245/heisenberg.pdf. Record: https://cvgmt.sns.it/paper/6245/. EMS Press, article record for the same title, DOI:10.4171/JEMS/1705. https://ems.press/journals/jems/articles/14299138. Journal text not compared with the bound manuscript. D. Jerison and J. M. Lee, Extremals for the Sobolev inequality on the Heisenberg group and the CR Yamabe problem, J. Amer. Math. Soc. 1 (1988), 1–13, DOI:10.1090/S0894-0347-1988-0924699-9. Cited as an external dependency; not independently re-proved in full.\par
\end{document}
